\section{Gaussian conditioning on the sphere}
\label{sec:spherical}

Conditioning independent centered Gaussian coordinates on their squared radius
gives an exact representation of a quadratic tilt of uniform sphere
measure. We quantify that representation in two stages. For bounded
deterministic spectra, radial-density estimates compare partition
integrals on the full sphere and its intersections with hyperplanes
through the origin with the same variational value. We then impose the
edge estimates of Section~\ref{sec:spectral} to bound small coordinate
projections and squared overlaps at criticality.

\subsection{Radial representation and partition estimates}

We begin with a deterministic ordered spectrum
$\lambda_1\ge\cdots\ge\lambda_n$, where $n\ge3$ and
$\max_i|\lambda_i|\le3$, and set
$\Lambda=\diag(\lambda_1,\ldots,\lambda_n)$.
Let $\sigma_n$ be uniform probability on
$\{u\in\R^n:\|u\|^2=n\}$.  For $0\le a\le2$, define
\begin{equation}
\label{eq:spherical-definitions}
 S(a)=\int e^{a u^T\Lambda u/2}\,\sigma_n(du),
 \qquad
 B(a)=\inf_{z>a\lambda_1}
 \left\{\frac{z-1}{2}
       -\frac1{2n}\sum_{i=1}^n\log(z-a\lambda_i)\right\}.
\end{equation}
The associated tilted probability is
\[
 \sigma_{\Lambda,a}(du)
 =S(a)^{-1}e^{a u^T\Lambda u/2}\,\sigma_n(du),
 \qquad
 \sigma_\Lambda=\sigma_{\Lambda,1}.
\]
Quadratic exponential laws on a sphere belong to the Bingham family
\cite{bingham1974}. Their representation through a Gaussian conditioned
on its radius, and the associated radial-density formula, are classical
in the analysis of Bingham and Fisher--Bingham normalizing constants
\cite{kumeWood2005,kumePrestonWood2013}. Gaussian integration and spherical
saddles also enter the SK analyses of Baik and Lee \cite{baikLee2016},
Landon \cite{landon2022}, and Du and Huang
\cite{duHuang2026FreeEnergy,duHuang2026Overlap}; the proof of Lemma~4.5 in
\cite{duHuang2026Overlap} uses the same representation by Gaussian
coordinates conditioned on the radius. We prove the exact identity
with the probability normalization and squared radius used here; no
saddlepoint approximation is used in this conditioning step.

\begin{lemma}[Gaussian representation]
\label{lem:spherical-identity}
For every $a\in[0,2]$, the infimum defining $B(a)$ has a unique minimizer
$z_a$.  With
\begin{equation}
\label{eq:spherical-saddle}
 k_i(a)=(z_a-a\lambda_i)^{-1},
 \qquad
 \sum_{i=1}^n k_i(a)=n,
\end{equation}
one has
\[
 1+a\lambda_n\le z_a\le1+a\lambda_1,
 \qquad
 \frac1{13}\le k_i(a)\le n.
\]
Let $Y_i$ be independent $N(0,k_i(a))$ variables, and let $p_a$ be the
density of $\sum_iY_i^2$.  Then
\begin{equation}
\label{eq:spherical-identity}
 S(a)=e^{nB(a)}
       \frac{p_a(n)}{p_{\chi_n^2}(n)}.
\end{equation}
The conditional law of $Y$ given $\|Y\|^2=n$, defined by radial
disintegration, is $\sigma_{\Lambda,a}$.
\end{lemma}

\begin{proof}
First choose the Gaussian variances so that the squared radius has
mean $n$. Then compare its radial density with the standard Gaussian
density at that same radius. To implement the first step, denote the
objective in \eqref{eq:spherical-definitions} by $F_a(z)$. On
$z>a\lambda_1$,
\[
 F_a'(z)=\frac12-\frac1{2n}\sum_i(z-a\lambda_i)^{-1},
 \qquad
 F_a''(z)=\frac1{2n}\sum_i(z-a\lambda_i)^{-2}>0.
\]
The derivative tends to $-\infty$ at the left endpoint and to $1/2$
at infinity.  The function itself tends to $+\infty$ at both ends.
This proves existence, uniqueness, and \eqref{eq:spherical-saddle}.

The average of the reciprocals of $z_a-a\lambda_i$ is one.  Their
minimum is therefore at most one, and their maximum is at least one,
giving the stated bounds on $z_a$.  Consequently
\[
 0<z_a-a\lambda_i
 \le1+a(\lambda_1-\lambda_n)\le13.
\]
The lower bound for $k_i(a)$ follows; the upper bound follows from
their positivity and sum $n$.  At $a=0$ these formulas give
$z_0=1$, $k_i(0)=1$, and $B(0)=0$.

The minimizer is now fixed. The radial comparison itself is exact
at any admissible $z$: put $Q=zI-a\Lambda$ and let
$Y^{z,a}\sim N(0,Q^{-1})$.  Its density relative to a standard
Gaussian vector at $y\in\R^n$ is
\[
 (\det Q)^{1/2}
 \exp\left(-\frac{z-1}{2}\|y\|^2+\frac a2y^T\Lambda y\right).
\]
The standard Gaussian has uniform conditional angular law.  Integrating
this ratio over the sphere of squared radius $n$ gives
\begin{equation}
\label{eq:spherical-radial-ratio}
 \frac{p_{z,a}(n)}{p_{\chi_n^2}(n)}
 =(\det Q)^{1/2}e^{-n(z-1)/2}S(a),
\end{equation}
where $p_{z,a}$ is the density of $\|Y^{z,a}\|^2$.
This identity is a polar-coordinate identity between continuous radial
densities at a positive radius.  At $z=z_a$, its rearrangement is
\eqref{eq:spherical-identity}.  The same density ratio, with the radius
fixed, proves the assertion about the conditional law.
\end{proof}

We will repeatedly use the two-square smoothing calculation also appearing in
\cite[Eq.~(7)]{BobkovNaumovUlyanov2021}.  If $G,H$ are
independent standard Gaussians and $d,e>0$, then the density
$r_{d,e}$ of $dG^2+eH^2$ satisfies
\begin{equation}
\label{eq:two-square-density}
 \begin{split}
 r_{d,e}(t)
 &=
 \frac1{2\pi\sqrt{de}}
 \int_0^t
 \frac{e^{-x/(2d)-(t-x)/(2e)}}{\sqrt{x(t-x)}}\,dx\\
 &\le \frac1{2\sqrt{de}},
 \qquad t>0.
 \end{split}
\end{equation}
The integral of $[x(t-x)]^{-1/2}$ over $(0,t)$ is $\pi$.
Convolving with further independent random variables preserves this
upper bound.

The next estimate is required at the mean of the weighted sum, including
arrays with a dominant weight. The Fourier majorant in its proof is the
weighted-square estimate of Bobkov, Naumov, and Ulyanov
\cite[Lemma~3]{BobkovNaumovUlyanov2021}, which we reproduce.
The compactness argument has the classical Gaussian-plus-quadratic
limiting structure described by Nourdin and Poly
\cite[Theorem~3.1]{NourdinPoly2012SecondChaos}; we give the needed
positive-weight argument directly. Their erratum
\cite{NourdinPoly2012Erratum} concerns separate finite-cumulant statements.

\begin{lemma}[Density at the mean]
\label{lem:mean-density}
There is a universal constant $c>0$ such that, for every $n\ge3$ and
every $k_1,\ldots,k_n>0$ with $\sum_i k_i=n$, the density $p$ of
$\sum_i k_iG_i^2$, for independent standard Gaussians $G_i$, satisfies
\begin{equation}
\label{eq:mean-density}
 p(n)\ge c\left(\sum_i k_i^2\right)^{-1/2}.
\end{equation}
\end{lemma}

\begin{proof}
After normalization, it is enough to exclude a sequence of densities
vanishing at zero. There are two ways such a sequence of weights can
behave: all weights vanish individually, giving a Gaussian limit, or a
positive largest weight remains and supplies a density bounded below
near zero. The proof treats these cases separately.

Set $\tau=(\sum_i k_i^2)^{1/2}$ and $d_i=k_i/\tau$. The centered
random variable
\[
 T_d=\sum_i d_i(G_i^2-1)
     =\frac{\sum_i k_iG_i^2-n}{\tau}
\]
has density $f_d$ with $f_d(0)=\tau p(n)$, and
$\sum_i d_i^2=1$.  Thus it suffices to bound $f_d(0)$ below uniformly
over all such finite arrays.  Each array has at least three positive
weights.  Its characteristic function is integrable, since its
absolute value decays as $O(|t|^{-n/2})$ at infinity; we always take
the resulting continuous version of its density.

Suppose, to the contrary, that arrays $d^{(r)}$ satisfy
$f_{d^{(r)}}(0)\to0$.  Sort each array in nonincreasing order and
extend it by zeros.  A diagonal subsequence has
$d_i^{(r)}\to c_i$ for every $i$, with
$\sum_i c_i^2\le1$.  We consider the two possible values of $c_1$.

\smallskip
\noindent\emph{Case $c_1=0$.}
For $d\ge0$, define
\[
 \psi_d(t)=\E e^{itd(G^2-1)}
          =e^{-itd}(1-2itd)^{-1/2}.
\]
For logarithms of characteristic functions, use the continuous
branch equal to zero at $t=0$.
For fixed $t$ and $|td|$ small, expansion gives
\[
 \log\psi_d(t)=-t^2d^2+O(|t|^3d^3).
\]
It follows that the characteristic function
$\phi_r(t)=\prod_i\psi_{d_i^{(r)}}(t)$ converges pointwise to
$e^{-t^2}$, because
$\sum_i(d_i^{(r)})^3\le d_1^{(r)}\sum_i(d_i^{(r)})^2=d_1^{(r)}$.

Weak convergence alone does not control a density at one point.
To pass through Fourier inversion we need an integrable bound uniform
in $r$. Concavity of $\log(1+x)$ and
$\sum_i(d_i^{(r)})^2=1$ give
\[
 |\phi_r(t)|
 =\prod_i\bigl(1+4(d_i^{(r)})^2t^2\bigr)^{-1/4}
 \le
 \bigl(1+4(d_1^{(r)})^2t^2\bigr)^{-1/(4(d_1^{(r)})^2)}.
\]
The function $x\mapsto \log(1+Ax)/x$ is decreasing for $A\ge0$.
Once $(d_1^{(r)})^2\le1/8$, the last expression is at most
$(1+t^2/2)^{-2}$.  Fourier inversion and dominated convergence imply
\[
 f_{d^{(r)}}(0)
 =\frac1{2\pi}\int_{\R}\phi_r(t)\,dt
 \longrightarrow
 \frac1{2\pi}\int_{\R}e^{-t^2}\,dt
 =\frac1{2\sqrt\pi},
\]
a contradiction.

\smallskip
\noindent\emph{Case $c_1>0$.}
Remove the largest coordinate and write
\[
 R_r=\sum_{i\ge2}d_i^{(r)}(G_i^2-1).
\]
We claim that $R_r$ converges in distribution to
\begin{equation}
\label{eq:density-residual-limit}
 R=\sum_{i\ge2}c_i(G_i^2-1)+Z,
 \qquad
 Z\sim N\left(0,\,2\left(1-\sum_{i\ge1}c_i^2\right)\right),
\end{equation}
where $Z$ and the displayed Gaussians are independent.
The series converges in $L^2$, since the tail beyond $m$ has variance
$2\sum_{i>m}c_i^2\to0$.  To identify the limit, note that the sorted weights
satisfy $d_{m+1}^{(r)}\le(m+1)^{-1/2}$.  For fixed $t$ and sufficiently
large $m$, the same logarithmic expansion as above gives, uniformly
in $r$,
\[
 \log\E e^{itR_r}
 =
 \sum_{i=2}^m\log\psi_{d_i^{(r)}}(t)
 -t^2\left(1-\sum_{i=1}^m(d_i^{(r)})^2\right)
 +O\left(\frac{|t|^3}{\sqrt{m+1}}\right).
\]
First let $r\to\infty$ and then $m\to\infty$.  The resulting
characteristic function is exactly that of \eqref{eq:density-residual-limit},
which proves the claim.

The residual limit has positive mass arbitrarily close to zero.
This is the fact that allows the surviving largest coordinate to
control the convolution density. For any $\delta>0$, choose a finite
truncation of its series
whose remaining tail variance is less than $\delta^2/16$.
Each nonzero summand in the finite part has a density positive near
zero, so the finite part together with $Z$ has positive probability
in $(-\delta/2,\delta/2)$.  Zero summands or a degenerate $Z$ cause
no difficulty.  The independent tail has positive probability in
the same interval by Chebyshev's inequality.  Intersecting these
events proves the assertion.  This argument uses square summability
of $(c_i)$, not summability.

In particular, by convergence in distribution and the open-set
inequality,
\[
 \liminf_{r\to\infty}\Prob\{|R_r|<c_1/3\}
 \ge\Prob\{|R|<c_1/3\}>0.
\]
The density of $d(G^2-1)$, for $d>0$, is
\[
 g_d(x)=
 \frac{e^{-(x+d)/(2d)}}{\sqrt{2\pi d(x+d)}}
 \,\mathbf1_{\{x>-d\}}.
\]
For $d$ in a sufficiently small neighborhood of $c_1$, this density
is bounded below by a positive constant depending on $c_1$ whenever
$|x|\le c_1/3$.  This interval stays strictly inside its support.

Independence now yields the pointwise convolution identity
\[
 f_{d^{(r)}}(0)=\E\bigl[g_{d_1^{(r)}}(-R_r)\bigr].
\]
To justify it at this particular point, rather than only almost
everywhere, the residual $R_r$ has a bounded density by
\eqref{eq:two-square-density}, using two of its positive coordinates.
Convolution of that bounded density with the $L^1$ function
$g_{d_1^{(r)}}$ is continuous: its change under translation by $h$
is bounded by the residual density supremum times
$\|g_{d_1^{(r)}}(\,\cdot+h)-g_{d_1^{(r)}}\|_{L^1}$, which tends to zero.
It therefore equals the continuous density $f_{d^{(r)}}$ everywhere.
Consequently
\[
 \liminf_{r\to\infty}f_{d^{(r)}}(0)
 \ge c(c_1)\liminf_{r\to\infty}\Prob\{|R_r|<c_1/3\}>0,
\]
again a contradiction.  The constant may depend on the subsequential
limit: any positive liminf contradicts the chosen sequence tending
to zero.  The two cases prove a universal lower bound and hence
\eqref{eq:mean-density}.
\end{proof}

We now use the radial estimates to compare $S(a)$ with $e^{nB(a)}$.
The orientation argument also fixes one sphere vector and integrates
a second over its orthogonal hyperplane. For that step, we need the same
upper comparison on every such hyperplane.

\begin{lemma}[Partition bounds on the sphere and hyperplanes]
\label{lem:partition-bounds}
There are universal constants $c,C>0$ such that, for every $n\ge3$,
every ordered spectrum $\lambda_1\ge\cdots\ge\lambda_n$ satisfying
$\max_i|\lambda_i|\le3$, and every $a\in[0,2]$,
\[
 c n^{-1/2}e^{nB(a)}
 \le S(a)\le C n^{1/2}e^{nB(a)}.
\]
If $H\subset\R^n$ is any subspace of dimension $n-1$, and
$\sigma_{n,H}$ is uniform probability on
$\{u\in H:\|u\|^2=n\}$, then
\[
 S_H(a):=\int e^{a u^T\Lambda u/2}\,\sigma_{n,H}(du)
 \le C n^{1/2}e^{nB(a)}.
\]
The same constant $C$ works for all $H$ and all $a\in[0,2]$.
\end{lemma}

\begin{proof}
For the full sphere, the radial identity and the density-at-the-mean
lemma give both bounds. For a hyperplane, only an upper density bound
is needed; interlacing allows us to use the full-space saddle without
asserting it is a saddle for the compression.

The chi-square density is
\[
 p_{\chi_m^2}(r)
 =\frac{r^{m/2-1}e^{-r/2}}{2^{m/2}\Gamma(m/2)},
 \qquad r>0.
\]
Stirling's formula gives
\[
 c n^{-1/2}\le p_{\chi_m^2}(n)\le C n^{-1/2}
 \quad\text{for }m=n\text{ or }m=n-1,\quad n\ge3;
\]
the constants absorb the finitely many small dimensions.
By Lemma~\ref{lem:mean-density} and
$\bigl(\sum_i k_i(a)^2\bigr)^{1/2}\le\sum_i k_i(a)=n$,
we have $p_a(n)\ge c/n$.
By \eqref{eq:two-square-density} and $k_i(a)\ge1/13$,
we also have $\sup_r p_a(r)\le C$.
Lemma~\ref{lem:spherical-identity} proves the two full-space bounds.

We now pass to a hyperplane. Let $\Lambda_H$ denote the compression
of $\Lambda$ to $H$, and write
its eigenvalues as $\theta_1\ge\cdots\ge\theta_{n-1}$.
The min--max principle gives
$\lambda_i\ge\theta_i\ge\lambda_{i+1}$.
At the full-space minimizer $z=z_a$, the matrix
$Q_H=z_a I_H-a\Lambda_H$ is positive definite.  Since $a\ge0$,
interlacing gives
\begin{equation}
\label{eq:spherical-compression-determinant}
 \begin{split}
 (\det Q_H)^{-1/2}
 &\le\prod_{i=1}^{n-1}(z_a-a\lambda_i)^{-1/2}\\
 &=\left(\prod_{i=1}^n(z_a-a\lambda_i)^{-1/2}\right)
   \sqrt{z_a-a\lambda_n}.
 \end{split}
\end{equation}
The last square root is at most $\sqrt{13}$.

Let $p_{H,a}$ be the squared-norm density of a Gaussian on $H$ with
covariance $Q_H^{-1}$.  The same polar-coordinate calculation as in
\eqref{eq:spherical-radial-ratio}, now in dimension $n-1$ but still
at squared radius $n$, yields
\[
 S_H(a)
 =e^{n(z_a-1)/2}(\det Q_H)^{-1/2}
   \frac{p_{H,a}(n)}{p_{\chi_{n-1}^2}(n)}.
\]
Every eigenvalue of $Q_H$ is at most
$z_a-a\lambda_n\le13$, so all its Gaussian variances are at least
$1/13$.  There are at least two of them, and
\eqref{eq:two-square-density} gives $p_{H,a}(n)\le C$.
Combining this bound with
\eqref{eq:spherical-compression-determinant} proves the assertion.
Only an upper density bound is used: the compressed Gaussian need
not have mean squared radius $n$, and $z_a$ need not minimize its
own variational problem.
\end{proof}

\subsection{Critical projections and overlap}

The preceding bounds require only a bounded deterministic spectrum.
We now impose Proposition~\ref{prop:spectral} to obtain the critical
projection and squared-overlap estimates used in Section~\ref{sec:orientation}.
Fix any deterministic $L_n\to\infty$ with $L_n\le\log n$, put
$s_0=L_n n^{-2/3}$, and take spectra in that proposition's good event.
All statements below hold for all sufficiently large $n$, uniformly
over these spectra.  At $a=1$, abbreviate
\[
 k_i=k_i(1),\qquad
 Y_i\sim N(0,k_i)\ \text{independently},\qquad
 p=p_1.
\]
For each $i$, let $p_{-i}$ be the density of $\sum_{j\ne i}Y_j^2$.

\begin{lemma}[Gaussian parameters at criticality]
\label{lem:critical-sphere}
With the preceding notation,
\begin{equation}
\label{eq:critical-spherical-parameters}
 \lambda_1<z_1<2+s_0,\qquad
 \left(\sum_i k_i^2\right)^{1/2}\le Cn\sqrt{s_0},
 \qquad
 k_j\ge c/s_0\quad(1\le j\le3).
\end{equation}
Moreover,
\[
 p(n)\ge \frac{c}{n\sqrt{s_0}},
 \qquad
 \sup_{r\in\R}p_{-i}(r)\le Cs_0
 \quad(1\le i\le n).
\]
\end{lemma}

\begin{proof}
We compare the exact saddle with the explicit point $2+s_0$. The first
trace locates the saddle and controls the total increase in variances
when moving to it; the second trace bounds the variances at the explicit
point. Set $h_i=(2+s_0-\lambda_i)^{-1}$,
whose denominators are positive by the edge estimates.  Since
$R(s)=1-\sqrt{s}+O(s)$ as $s\downarrow0$,
Proposition~\ref{prop:spectral} gives, uniformly over the good spectra,
\[
 \sum_i h_i
 =n\bigl(1-\sqrt{s_0}+o(\sqrt{s_0})\bigr)<n.
\]
The function $z\mapsto\sum_i(z-\lambda_i)^{-1}$ is strictly decreasing
from infinity on $(\lambda_1,\infty)$.  Its value at $z_1$ is $n$;
hence $\lambda_1<z_1<2+s_0$ and $k_i\ge h_i$.
Because $\sum_i k_i=n$, the total increase is exactly
\[
 \sum_i(k_i-h_i)
 =n-\sum_i h_i
 =n\sqrt{s_0}\bigl(1+o(1)\bigr).
\]
A nonnegative vector has Euclidean norm at most the sum of its coordinates.
The triangle inequality and the second resolvent bound therefore give
\[
 \begin{split}
 \left(\sum_i k_i^2\right)^{1/2}
 &\le \left(\sum_i h_i^2\right)^{1/2}
       +\sum_i(k_i-h_i)\\
 &\le C\sqrt n\,s_0^{-1/4}+Cn\sqrt{s_0}
 \le Cn\sqrt{s_0}.
 \end{split}
\]
For the last step,
$n^{-1/2}s_0^{-3/4}=L_n^{-3/4}\to0$.
For $j\le3$, the edge estimates also give
\[
 0<z_1-\lambda_j\le2+s_0-\lambda_j\le Cs_0,
\]
which proves the lower bounds on the first three variances.
Lemma~\ref{lem:mean-density} now gives the asserted lower bound on
$p(n)$.

The three leading variances are needed for a uniform deletion
estimate: after deleting any one index $i$, at least two of the
indices $1,2,3$ remain.  Applying \eqref{eq:two-square-density} to their squares,
and then convolving with the other squares, yields
$\sup_r p_{-i}(r)\le Cs_0$.
\end{proof}

\begin{proposition}[Small coordinate projections on the tilted sphere]
\label{prop:spherical-small-ball}
For every $b>0$, the preceding good spectra satisfy
\begin{equation}
\label{eq:spherical-small-ball}
 \sigma_\Lambda\{|u_1|\le b n^{1/3}\}
 \le C L_n^2 b.
\end{equation}
In particular, given any positive deterministic $b_n\to0$, one can
choose $L_n\to\infty$ with $L_n\le\log n$ and $L_n^2b_n\to0$, so that
the probability in \eqref{eq:spherical-small-ball} tends to zero
uniformly over the resulting good spectra.
\end{proposition}

\begin{proof}
Conditioning a single coordinate makes the role of the radial
estimates transparent. Its unconditioned Gaussian mass is small, and
the ratio of the residual radial density to the full radial density
quantifies the cost of conditioning. Let $\varphi_k$ be the density of
$N(0,k)$. Lemma~\ref{lem:spherical-identity} gives,
for $T>0$,
\[
 \begin{split}
 \sigma_\Lambda\{|u_1|\le T\}
 &=\frac1{p(n)}
   \int_{|y|\le T}\varphi_{k_1}(y)p_{-1}(n-y^2)\,dy\\
 &\le \frac{Cs_0}{p(n)}\,\Prob\{|Y_1|\le T\}.
 \end{split}
\]
The densities of the nonnegative residual sums are taken to be zero
at negative arguments.  By Lemma~\ref{lem:critical-sphere} and the
Gaussian density bound,
\[
 \frac{Cs_0}{p(n)}\le Cn s_0^{3/2},
 \qquad
 \Prob\{|Y_1|\le b n^{1/3}\}
 \le \frac{Cb n^{1/3}}{\sqrt{k_1}}
 \le Cb n^{1/3}\sqrt{s_0}
 =Cb\sqrt{L_n}.
\]
Their product is $CbL_n^2$, proving
\eqref{eq:spherical-small-ball}.  For the final assertion, for example,
$L_n=\min\{\log n,b_n^{-1/4}\}$ for all sufficiently large $n$ has
the required properties.
\end{proof}

\begin{lemma}[Coordinate moments and squared overlap]
\label{lem:spherical-covariance}
The probability $\sigma_\Lambda$ has mean zero and diagonal
covariance in the eigenbasis of $\Lambda$.  For every coordinate,
\begin{equation}
\label{eq:spherical-coordinate-moment}
 \int u_i^2\,\sigma_\Lambda(du)
 \le \frac{Cs_0 k_i}{p(n)}.
\end{equation}
Consequently, for independent $U,V$ with law $\sigma_\Lambda$,
\begin{equation}
\label{eq:spherical-overlap}
 \E\left(\frac{U^TV}{n}\right)^2
 \le
 \frac{Cs_0^2}{n^2p(n)^2}\sum_i k_i^2
 \le Cn^{-2/3}L_n^4.
\end{equation}
\end{lemma}

\begin{proof}
We reuse the same one-coordinate conditioning, now with the weight
$y^2$. First, the defining density of $\sigma_\Lambda$ is unchanged by
changing the sign of any individual coordinate.  This proves the assertions
about its mean and off-diagonal covariances.
Conditioning as in the preceding proof, now with the nonnegative
integrand $y^2$, gives
\[
 \int u_i^2\,\sigma_\Lambda(du)
 =\frac1{p(n)}
   \int_{\R}y^2\varphi_{k_i}(y)p_{-i}(n-y^2)\,dy
 \le\frac{Cs_0}{p(n)}\,\E Y_i^2
 =\frac{Cs_0 k_i}{p(n)}.
\]
Write $m_i=\int u_i^2\,\sigma_\Lambda(du)$.
Independence of $U,V$ and diagonality give
\[
 \E\left(\frac{U^TV}{n}\right)^2
 =\frac1{n^2}\sum_i m_i^2
 \le\frac{Cs_0^2}{n^2p(n)^2}\sum_i k_i^2.
\]
Lemma~\ref{lem:mean-density} bounds $p(n)^{-2}$ by
$C\sum_i k_i^2$, and Lemma~\ref{lem:critical-sphere} bounds the latter
sum by $Cn^2s_0$.  The resulting upper bound is
$Cn^2s_0^4=Cn^{-2/3}L_n^4$, as claimed.
\end{proof}
